% !TeX root = ../thuthesis-example.tex

\chapter{总结与展望}

\section{全文总结}

本文围绕 LingoBridge 国际中文学习平台，完成了基于 DKT 知识追踪与 Levenshtein 距离校准的国际中文自适应学习推荐系统设计与原型实现。系统以真实教学业务流程为基础，将课程、作业、录音练习、批改反馈、学习时长和班级排名等数据转化为可分析的学习行为记录，并通过数据清洗、特征构建和指标聚合形成学生学习画像。

在算法层面，系统采用 DKT 模型预测学生对拼音、声调和词汇知识点的动态掌握状态，使用 Levenshtein 编辑距离对 ASR 字幕或拼音转写偏误进行校准，并结合知识点依赖关系生成可解释练习推荐。该方案避免了本地演示中过度依赖外部大模型服务的问题，同时能够清楚体现课程要求中的数据预处理、算法设计、编码实现、运行测试和报告撰写等环节。

在工程层面，项目保留 LingoBridge 已有 MVP 主流程，将 AI 模块作为独立的 Python 分析子工程进行组织，便于后续对接学生端首页、个人排名页、教师学生总览页和管理端数据中台。整体方案兼顾了课程设计可展示性、算法可解释性和项目已有业务场景的连续性。

\section{系统运行效果说明}

当前原型已经能够完成学习数据生成、知识点依赖分析、DKT 模型训练、Levenshtein 字幕校准和推荐结果导出。报告中的实验章节展示了知识图谱拓扑分析、模型训练曲线、ROC 对比、薄弱点推荐案例和字幕校准对照，能够支撑课堂答辩中对“人工智能方法如何应用到系统”的说明。

后续在应用演示时，可将分析结果进一步接入三端界面：学生端展示个人掌握度、薄弱知识点和下一步练习建议；教师端展示班级风险学生、共性错误和作业反馈建议；管理端展示全局学习指标、模块运行状态和平台级数据分析看板。

\section{研究不足与未来展望}

本项目目前仍以大学学校单租户为主要目标，数据规模以模拟数据和 MVP 业务结构为基础，尚未达到真实长期教学场景中的大规模、多学期、多班级数据量。因此，DKT 模型在小样本条件下的效果可能不稳定，部分场景中规则基线仍具有更好的鲁棒性。未来可以在真实学习日志积累后重新训练模型，并尝试引入更适合小样本教育数据的序列建模方法。

在文本校准方面，Levenshtein 距离具备实现简单、可解释性强和本地运行稳定的优势，但它主要刻画字符串层面的相似度，无法充分理解上下文语义。未来可以在保持可解释基线的基础上，引入拼音混淆矩阵、声调错误模型或轻量检索增强方法，提高对真实 ASR 偏误的处理能力。

在系统演进方面，后续工作应重点完成 AI 模块与前端页面的数据对接，使推荐结果真正进入学生学习路径和教师教学决策流程。同时，需要持续完善数据隐私、日志留痕、模型版本管理和本地演示进程管理，保证系统在课程答辩和后续开发中都具备稳定、可回滚、可解释的工程质量。
